\begin{proof}[Proof of \cref{obj:lem:exact-covariance}]
\begin{enumerate}
\item We first establish the variance bound and choose its constant. Applying \cref{obj:lem:positive-endpoint} with \(\beta=1\), choose an absolute constant \(C_{\mathrm{end}}>0\) such that
\[
\int_0^1 K_L(x)^2\,dx\le C_{\mathrm{end}}L^2
\qquad(L\ge1).
\]
Define
\[
C=\max\{C_{\mathrm{end}},1296\}.
\]
Then \(C>0\). Fix \(L\ge1\) and \(\sigma\in[0,1]\). We have
\[
m_L+1=3L-2\le3L,
\qquad
\int_0^1 K_L(x)^2\,dx\le CL^2.
\]
The numerator \(1-T_L(1-2x)\) has degree at most \(L\); its polynomial quotient by \(2x\) has degree at most \(L-1\). Cubing and multiplying by the normalization constant in \cref{obj:def:endpoint-kernel} therefore gives
\[
\deg K_L\le3(L-1)=m_L.
\]

For \(1\le j\le m_L\), squaring the nonnegative sides of \cref{obj:lem:factorial-derivatives}, applied to \(K_L\), gives
\[
\int_0^1 [K_L^{(j)}(x)]^2\,dx
\le
\frac{16^j(m_L+1)^{4j}}{(j!)^2}
\int_0^1 K_L(x)^2\,dx.
\]
For \(j=0\), this inequality is an equality. Thus it holds throughout \(0\le j\le m_L\), including when \(L=1\). Moreover,
\[
16^j(m_L+1)^{4j}
=\bigl(16(m_L+1)^4\bigr)^j
\le(1296L^4)^j
\le(CL^4)^j.
\]
Consequently, for each such \(j\),
\[
\begin{aligned}
\frac{\sigma^{2j}}{j!}
\int_0^1 [K_L^{(j)}(x)]^2\,dx
&\le
\frac{\sigma^{2j}}{j!}
\frac{16^j(m_L+1)^{4j}CL^2}{(j!)^2}\\
&\le
CL^2\frac{(C\sigma^2L^4)^j}{(j!)^3}.
\end{aligned}
\]
Every summand on the left is nonnegative. Summing and using \(3/4\le1\) proves
\[
V_L(\sigma)
\le
\sum_{j=0}^{m_L}\frac{\sigma^{2j}}{j!}
\int_0^1 [K_L^{(j)}(x)]^2\,dx
\le
CL^2\sum_{j=0}^{m_L}
\frac{(C\sigma^2L^4)^j}{(j!)^3}.
\]
% lean: CausalSmith.Stat.RdTruesideNoiseFrontier.kernelVariance_factorial_envelope

\item Suppose first that \(\sigma=0\). In \cref{obj:def:inverse-heat}, the term \(k=0\) equals \(K_L\), and every term with \(k>0\) vanishes. Hence
\[
Q_{L,0}(w)=K_L(w).
\]
Since \(m_L\ge0\), the derivative sum contains its zeroth term. With the usual convention \(0^0=1\), its remaining terms vanish:
\[
\sum_{j=0}^{m_L}\frac{0^{2j}[K_L^{(j)}(x)]^2}{j!}
=K_L(x)^2.
\]
The two Gaussian expectations are therefore
\[
\mathbb E[Q_{L,0}(x+0\varepsilon)]=K_L(x),
\qquad
\mathbb E[Q_{L,0}(x+0\varepsilon)^2]=K_L(x)^2,
\]
as required.
% lean: CausalSmith.Stat.RdTruesideNoiseFrontier.inverseHeat_zero_covariance

\item Assume now that \(\sigma\ne0\), and fix \(x\in\mathbb R\). For a nonnegative integer \(m_{\mathrm{cut}}\), a real number \(\sigma_{\mathrm{heat}}\), and a real polynomial \(p\), define the polynomial operator
\[
(\mathcal T_{m_{\mathrm{cut}},\sigma_{\mathrm{heat}}}p)(w)
=
\sum_{k=0}^{\lfloor m_{\mathrm{cut}}/2\rfloor}
\frac{(-\sigma_{\mathrm{heat}}^2/2)^k}{k!}\,
p^{(2k)}(w).
\]
Thus \cref{obj:def:inverse-heat} gives
\[
Q_{L,\sigma}=\mathcal T_{m_L,\sigma}K_L.
\]
Define the affine polynomial and its monomial coefficients by
\[
p_{\mathrm{aff}}(z)=K_L(x+\sigma z)
=\sum_{j=0}^{m_L}c_{\mathrm{aff},j}z^j
\qquad(z\in\mathbb R).
\]
Its degree is at most \(m_L\), by the degree bound established above.

Repeated differentiation of the affine substitution gives, for every nonnegative integer \(j\),
\[
p_{\mathrm{aff}}^{(j)}(z)
=\sigma^jK_L^{(j)}(x+\sigma z).
\]
Substituting this identity into the definition of the operator yields
\[
\begin{aligned}
(\mathcal T_{m_L,1}p_{\mathrm{aff}})(z)
&=
\sum_{k=0}^{\lfloor m_L/2\rfloor}
\frac{(-1/2)^k}{k!}\,
\sigma^{2k}K_L^{(2k)}(x+\sigma z)\\
&=
(\mathcal T_{m_L,\sigma}K_L)(x+\sigma z).
\end{aligned}
\]
Differentiation distributes over finite sums and commutes with scalar multiplication. Applying this to the monomial decomposition of \(p_{\mathrm{aff}}\), we obtain
\[
Q_{L,\sigma}(x+\sigma z)
=
\sum_{j=0}^{m_L}
c_{\mathrm{aff},j}\,
(\mathcal T_{m_L,1}(z^j))(z).
\]
Here the operator acts on the polynomial \(z\mapsto z^j\) before evaluation.
% lean: CausalSmith.Stat.RdTruesideNoiseFrontier.finiteInverseHeat_affine_eval

\item We identify the transformed monomials with the probabilists' Hermite polynomials \(H_j\) of \cref{obj:lem:classical-hermite-facts}. For \(0\le j\le m_L\), all derivatives of \(z^j\) of order greater than \(j\) vanish. Enlarging the cutoff from \(j\) to \(m_L\) therefore adds zero terms, and
\[
\mathcal T_{m_L,1}(z^j)
=
\mathcal T_{j,1}(z^j)
=
\sum_{k=0}^{\lfloor j/2\rfloor}
\frac{(-1/2)^k}{k!}\,
\frac{j!}{(j-2k)!}\,z^{j-2k}.
\]

To compare coefficients, define the odd double factorial by
\[
(2k-1)!!=\prod_{k'=1}^{k}(2k'-1)
\qquad(k\ge0),
\]
where the empty product equals \(1\), so that \((-1)!!=1\).
Separating the even and odd factors of \((2k)!\) gives
\[
(2k)!=2^k k!(2k-1)!!.
\]
Expanding the generating function in \cref{obj:lem:classical-hermite-facts} coefficientwise gives, for nonnegative integers \(i,j\),
\[
\text{coefficient of \(z^i\) in \(H_j(z)\)}
=
\begin{cases}
(-1)^k\binom ji(2k-1)!!,
& j=i+2k,\quad k\ge0,\\
0,
& i>j\ \text{or }j-i\text{ is odd}.
\end{cases}
\]
Indeed, the coefficient of \(z^it^{i+2k}\) in
\(\exp(zt)\exp(-t^2/2)\) is
\((-1/2)^k/(i!k!)\); multiplying by \((i+2k)!\) and using the factorial identity above gives the displayed coefficient.

When \(j=i+2k\), this coefficient equals
\[
(-1)^k\binom ji(2k-1)!!
=
\frac{(-1/2)^k}{k!}\frac{j!}{i!}.
\]
This is exactly the coefficient of \(z^i\) in
\(\mathcal T_{j,1}(z^j)\). For \(i>j\), or for \(j-i\) odd, both coefficients are zero. Thus
\[
\mathcal T_{m_L,1}(z^j)=H_j(z)
\qquad(0\le j\le m_L).
\]
% lean: CausalSmith.Stat.RdTruesideNoiseFrontier.finiteInverseHeat_unit_monomial

Finally, finite Taylor expansion followed by scaling gives
\[
K_L(x+\sigma z)
=
\sum_{j=0}^{m_L}
\frac{\sigma^jK_L^{(j)}(x)}{j!}\,z^j,
\qquad
c_{\mathrm{aff},j}
=
\frac{\sigma^jK_L^{(j)}(x)}{j!}.
\]
Substituting these coefficients and the transformed monomials into the preceding decomposition proves the finite identity
\[
Q_{L,\sigma}(x+\sigma z)
=
\sum_{j=0}^{m_L}
\frac{\sigma^jK_L^{(j)}(x)}{j!}\,H_j(z)
\qquad(z\in\mathbb R).
\]
% lean: CausalSmith.Stat.RdTruesideNoiseFrontier.finiteInverseHeat_hermite_expansion

\item By \cref{obj:lem:classical-hermite-facts}, \(H_0=1\), and orthogonality with this constant polynomial gives
\[
\mathbb E[H_j(\varepsilon)]
=
\mathbb E[H_j(\varepsilon)H_0(\varepsilon)]
=
\begin{cases}
1,&j=0,\\
0,&j\ge1.
\end{cases}
\]
The Hermite polynomials are square integrable under the Gaussian law, so expectations pass through the finite sum. Its zeroth term is present because \(m_L+1>0\). Hence
\[
\mathbb E[Q_{L,\sigma}(x+\sigma\varepsilon)]
=
c_{\mathrm{aff},0}
=
K_L(x).
\]
% lean: CausalSmith.Stat.RdTruesideNoiseFrontier.hermite_finite_sum_mean

\item Expanding the square of the same finite sum and applying \cref{obj:lem:classical-hermite-facts} gives
\[
\begin{aligned}
\mathbb E[Q_{L,\sigma}(x+\sigma\varepsilon)^2]
&=
\sum_{j=0}^{m_L}\sum_{k=0}^{m_L}
c_{\mathrm{aff},j}c_{\mathrm{aff},k}
\mathbb E[H_j(\varepsilon)H_k(\varepsilon)]\\
&=
\sum_{j=0}^{m_L}c_{\mathrm{aff},j}^{\,2}j!\\
&=
\sum_{j=0}^{m_L}
\frac{\sigma^{2j}[K_L^{(j)}(x)]^2}{j!}.
\end{aligned}
\]
The products are integrable by square integrability and the Cauchy--Schwarz inequality, which justifies the finite expansion under expectation. The last equality uses \(j!>0\) for every \(j\ge0\). Together with the zero-noise case and the variance bound, this proves the result with the chosen absolute constant \(C\).
% lean: CausalSmith.Stat.RdTruesideNoiseFrontier.hermite_finite_sum_square
\end{enumerate}
\end{proof}
